\documentclass[10pt,a4paper]{amsart}
\usepackage[margin=19mm]{geometry}
\usepackage{amsmath,amssymb,mathtools}
\usepackage[scr=rsfs,cal=euler]{mathalfa}
\usepackage{xfrac}
\usepackage[unicode,hidelinks,bookmarks=false]{hyperref}
\newcommand{\ie}{i.e.\ }
\newcommand{\cf}{cf.\ }
\newcommand{\resp}{respectively\ }
\newcommand{\Subsection}{\subsection}
\providecommand{\nobreakdash}{\nobreak\hbox{-}}
\setlength{\emergencystretch}{2em}
\multlinegap0pt
\usepackage{bm}
\usepackage{bbm}
\usepackage{dsfont}
\usepackage{xspace}
\usepackage{cancel}
\usepackage{mathtools}
\usepackage{amsthm}
\makeatletter
\@ifundefined{thm}{\newtheorem{thm}{Thm}}{}
\@ifundefined{lemma}{\newtheorem{lemma}[thm]{Lemma}}{}
\makeatother
\newcommand{\Com}{{\mathsf{ComRings}}}
\newcommand{\Ker}{{\rm Ker  }}
\newcommand{\N}{{\mathbb N}}
\newcommand{\Z}{{\mathbb Z}}
\newcommand{\ab}[1]{\langle {#1} \rangle}
\newcommand{\ac}[1]{\{#1\}}
\newcommand{\ilim}{\mathop{\varprojlim}\limits} % inverse limit
\renewcommand{\tilde}{\widetilde}
\title[A universal group-theoretic characterisation of $p$-typical ]{A universal group-theoretic characterisation of $p$-typical Witt vectors}
\author{Supriya Pisolkar, Biswanath Samanta}
\date{Source: arXiv:2405.12680v2 (CC BY 4.0)}
\begin{document}
\maketitle

\noindent\emph{Source and licence.} A universal group-theoretic characterisation of $p$-typical Witt vectors. Version arXiv:2405.12680v2 (CC BY 4.0), distributed under the Creative Commons Attribution 4.0 International licence, as recorded in the arXiv licence metadata for this exact version and shown on the abstract page https://arxiv.org/abs/2405.12680v2. Full rights notice: licenses/arxiv-2405.12680-CC-BY-4.0.txt.\par\smallskip

\noindent\textit{Editorial scope.} Counted unit: the Lemma whose source label is \texttt{genrelxa} in the article (source0.tex lines 450--452, source label \texttt{genrelxa}), delivered together with the article's own complete original proof of it (source0.tex lines 453--499, one \texttt{proof} environment of 47 lines). No mathematics is written, repaired or re-derived: statement and proof are the article's own text line for line. Required context retained uncounted: ilimxa (lines 297--306); xisprewitt (lines 313--314); ceqg (lines 372--374); zindcomm (lines 443--445). Every definition, display and label of those lines is retained unchanged. Omitted from this package and not redistributed: 1--296, 307--312, 315--371, 375--442, 446--449, 500--545. Nothing inside a retained excerpt is omitted or altered: every retained body is delivered exactly as the article's lines are written, so no omission receipt of any kind accompanies this package. Citations: the retained excerpts carry no citation command at all, so no bibliography entry of the article is transcribed and no reference is redistributed. Reference adaptation: none was needed and none was made; every reference inside the delivered unit resolves inside it, so no unresolved reference or citation is produced. Printed slips kept literal: no printed word, symbol, hypothesis or number of the counted statement or of its proof was altered. Layout and class: the article's own class is \texttt{amsart} with packages including geometry, amssymb, amsmath, amsthm, amstext, amscd, latexsym, graphics, graphicx, bbm, caption, xcolor, hyperref, tikz; the wrapper is the lane's pinned \texttt{amsart} editorial wrapper at 10pt on A4, which restates the theorem-like environments the retained text needs and adds the article's own macro definitions that the retained text needs (\texttt{\textbackslash Com}, \texttt{\textbackslash Ker}, \texttt{\textbackslash N}, \texttt{\textbackslash Z}, \texttt{\textbackslash ab}, \texttt{\textbackslash ac}, \texttt{\textbackslash ilim}, \texttt{\textbackslash tilde}), with extra wrapper packages (bm, bbm, dsfont, xspace, cancel, mathtools). The delivered numbering is the unit's own. Assets: no retained span contains a figure environment, a graphics-inclusion command, a PDF-page inclusion, a table or a TikZ picture; no image file is copied into this package, the package declares no asset dependency and no image file is redistributed with it. Licence: version arXiv:2405.12680v2 of this article is distributed under the Creative Commons Attribution 4.0 International licence, as recorded in the arXiv licence metadata for this exact version and shown on the abstract page https://arxiv.org/abs/2405.12680v2. A full rights notice is delivered at licenses/arxiv-2405.12680-CC-BY-4.0.txt. Characters: every retained excerpt is pure ASCII, so the delivered engine, XeLaTeX with its default Latin Modern text font, reports no missing character.
\begin{lemma}\label{ilimxa} 
For $R \in \Com$, let $\tilde{X}(R)\subset X(R)$ be the subgroup generated by \\
$\ac{V^n\ab{r} \ | \ n\in \N_0, \ r\in R}$. Then 
\begin{enumerate}
\item The natural map $X(R) \to \ilim_n \frac{X(R)}{V^n(X(R))}$ is an isomorphism. 
\item The natural map $\frac{\tilde{X}(R)}{V^n(\tilde{X}(R))} \to \frac{X(R)}{V^n(X(R))}$ is an isomorphism.
\end{enumerate}
In particular we have a canonical isomorphism 
$X(R) \cong \ilim_n\frac{\tilde{X}(R)}{V^n(\tilde{X}(R))}$.
\end{lemma}

\begin{lemma}\label{xisprewitt} $X : \Com \to \mathsf{Ab}$ is a pre-Witt functor.
\end{lemma}

\begin{lemma}\label{ceqg}
If $A=\Z[S]$ is a free commutative algebra, then $C(A) \cong G(A)$. In particular, $C(A)$ is $p$-torsion free.
\end{lemma}

\begin{lemma}\label{zindcomm} Let $p\neq 2$.
Let $A = \Z[S]$ be a commutative polynomial ring over a set $S$. Let $\{f_i\}_{i=1}^r$ be a finite set of distinct non-zero elements of $A$. Further assume that $f_i\neq -f_j$ for any $i\neq j$. Then the subset $\{\ab{f_i}\}_{i=1}^r$ of $X(A)$ is $\Z$-linearly independent. 
\end{lemma}

\begin{lemma} \label{genrelxa} Let $A=\Z[S]$ be a free polynomial algebra. 
There exists a canonical isomorphism  $C(A) \cong X(A)$. 
\end{lemma}

\begin{proof}
We know that $X(A)$ is a pre-Witt functor by Lemma \ref{xisprewitt}. Since $C$ is a universal weak pre-Witt functor there exists a natural transformation $\eta: C \to X$. For a polynomial algebra $A$, we have the group homomorphism $\eta: C(A) \to X(A)$ which is given by $V^n_a \mapsto V^n\ab{a}$.  By Lemma \ref{ilimxa}(1) this map is surjective. By Lemma \ref{ceqg} we know that $G(A)\cong C(A)$, thus composing with the group homomorphism $\tilde{G}(A)/\tilde{H}(A) \to G(A)$, we get the map 
 
$$\overline{\eta}: \tilde{G}(A)/\tilde{H}(A) \to G(A) \to X(A)$$
The image of $\overline{\eta}$ lands in $\tilde{X}(A)$. Thus we have the group epimorphism 
$$\tilde{\eta}:  \tilde{G}(A)/\tilde{H}(A) \to \tilde{X}(A);  \  \ \ V^n_a \mapsto V^n\ab{a} .$$

\noindent We will now prove the theorem in following steps.\\

\noindent {\underline{Step 1}:} To prove the theorem it is enough to show that $\tilde{\eta}$ is an isomorphism: This is because, if $\tilde{\eta}$ is an isomorphism then since $\tilde{\eta}$ is compatible with $V$, after taking the $V$-completion we get that $$G^0(A) \cong X(A).$$ Since $X(A)$ has no $p$-torsion, this implies, using Lemma \ref{ceqg} $$C(A) \cong G(A) \cong X(A).$$

\noindent Thus, it is enough to show that the group epimorphism $\tilde{\eta}$ is injective.  This we will prove in the next step.\\

\noindent {\underline{Step 2}:} To show that $\tilde{\eta}$ is injective, we in fact consider the map (again denoted by $\tilde{\eta}$);
$$\tilde{\eta}: \tilde{G}(A) \to \tilde{X}(A)$$
We will show that $\Ker(\tilde{\eta}) = \tilde{H}(A)$. It is easy to observe from the definition of $\tilde{H}(A)$ that $\tilde{H}(A) \subseteq \Ker(\tilde{\eta})$. \\

Let  $\alpha= \sum_{i=1}^\ell c_i V^{n_i}_{a_i} \in \Ker(\tilde{\eta})$. We will prove that $\alpha\in \tilde{H}(A)$, i.e. 
$p^k \alpha \in H(A)$ for some $k \geq 0$. We will establish this in following two steps.\\

\noindent {\underline{Step 2a}:} We first prove the case when $n_i=0$ for all $i$, i.e. 
$$ \alpha = \sum_{i=1}^\ell c_iV^0_{a_i}$$ 
We may assume without loss of generality that $a_i$ are all distinct. Moreover if $a_r=-a_s$ for some $r\neq s$ then 
$$ V^0_{a_r} = - V^0_{a_s} \ {\rm  mod} \ H(A).$$ Thus we may replace the expression $c_rV^0_{a_r}+c_sV^0_{a_s}$ with  $(c_s-c_r)V^0_{a_s}$. In particular we may assume that not only $a_i$ are distinct, but $a_i\neq -a_j$ for any $i\neq j$. In this case, $c_i$ must be zero for all $i$ by Lemma \ref{zindcomm}. Thus $\alpha=0$ and hence is in $\tilde{H}(A)$. \\

\noindent {\underline{Step 2b}:} For $n_i \geq 0$, we prove the claim by induction on $n_{\alpha}: = {\rm max}_i\ac{n_i}- {\rm min}_i \ac{n_i}$. Since $\tilde{X}(A) \xrightarrow{V} \tilde{X}(A)$ and $\tilde{H}(A)\xrightarrow{V} \tilde{H}(A)$ are both injective maps, we may assume without loss of generality that ${\rm min}_i \ac{n_i} = 0$.  The starting step of induction is proved in Step 1. Now by rearranging, we may assume that $n_i = 0$ for $1\leq i \leq s$ and $n_i>0$ for $s < i\leq \ell$. Thus without loss of generality we have 
$$ \alpha = \sum_{i=1}^s c_iV^0_{a_i} + \sum_{i>s}^\ell c_iV^{n_i}_{a_i}$$

\noindent We need to show $p^k \alpha \in H(A)$, for some $k \geq 0$.
\begin{align*}
p\alpha & = \sum_{i=1}^s pc_iV^0_{a_i} + \sum_{i>s}^\ell pc_iV^{n_i}_{a_i} \\ 
       & = \sum_{i=1}^s c_i (pV^0_{a_i}-V^1_{a_i^p}) +\sum_{i=1}^s c_iV^1_{a_i^p} + \sum_{i>s}^\ell pc_iV^{n_i}_{a_i} \\
            & = \beta + \sum_{i=1}^s c_iV^1_{a_i^p} + \sum_{i>s}^\ell pc_iV^{n_i}_{a_i} \\
            & =: \beta + \gamma 
\end{align*}
Since $\alpha \in \Ker(\tilde{\eta})$ we have 
$$p\tilde{\eta}(\alpha) = \tilde{\eta}(\beta) + \tilde{\eta}(\gamma) = (p\sum_{i=1}^sc_ia_i, \cdots ) = (0, 0, 0, \cdots).$$
This implies that $\sum_{i=1}^s c_ia_i = 0$ since $X(A)$ is p-torsion free.\\

\noindent We now claim that $\beta \in H(A)$: To prove this claim we will use the defining relations of $H(A)$.\\
 Consider the map  $\phi: A \to \tilde{G}(A)$ given by $ a \mapsto V^1_{a^p}-pV^0_{a}$ which is a linear map, i.e. 
 $$\phi(\sum_{i=1}^s c_ia_i) \equiv \sum_{i=1}^sc_i \phi(a_i) (\text{mod} \  H(A))$$
 Thus, $\beta \equiv 0 (\text{mod} \  H(A))$. This proves that $\beta \in H(A)$ and hence is in $\tilde{H}(A)$. Thus, $\beta \in \Ker(\tilde{\eta})$. This implies that $\gamma \in \Ker(\tilde{\eta})$. Moreover observe that $n_{\gamma} < n_{\alpha}$, hence by the induction hypothesis, $\gamma \in \tilde{H}(A)$.  Thus, $p^{r}\gamma \in H(A)$ for some $r \geq 0$. Thus, $p^{r+1}\alpha \in H(A)$. 
 Thus, $\Ker(\tilde{\eta}) \subseteq \tilde{H}(A)$. \\ 
 
\noindent This proves that $\tilde{G}(A)/ \tilde{H}(A) \cong \tilde{X}(A)$ and hence the theorem. 
\end{proof}

\end{document}
